%------------------------------------------------------------------------------%

\usepackage{calculo_varias_variaveis-1}

\title{Derivadas Parciais de Ordem Superior}

%------------------------------------------------------------------------------%
\begin{document}

\addtitle

%------------------------------------------------------------------------------%
\section{Derivadas Parciais de Segunda Ordem}

%------------------------------------------------------------------------------%
\begin{frame}{Conceito}
  Derivada de uma derivada de uma função
\end{frame}

%------------------------------------------------------------------------------%
\begin{frame}{Notações}
\begin{columns}
\column{0.5\linewidth}
  Derivadas de segunda ordem
  \[
    \frac{\partial^2 f}{\partial x^2} \qquad\qquad\qquad
    f_{xx}
  \]

  \vspace{\baselineskip}
  \[
    \frac{\partial^2 f}{\partial y^2} \qquad\qquad\qquad
    f_{yy}
  \]
\column{0.5\linewidth}
  \pause
  Derivadas de segunda ordem mistas
  \[
    \frac{\partial^2 f}{\partial x \partial y} \qquad\qquad\qquad
    f_{yx}
  \]

  \vspace{\baselineskip}
  \[
    \frac{\partial^2 f}{\partial y \partial x} \qquad\qquad\qquad
    f_{xy}
  \]
\end{columns}
\end{frame}

%------------------------------------------------------------------------------%
\begin{frame}{Atenção Para a Ordem das Derivadas Mistas}
  \large
  \[
    \frac{\partial^2 f}{{\color{magenta} \partial x} {\color{blue} \partial y}}
    = \frac{\partial}{{\color{magenta} \partial x}}
      \left(\,\frac{\partial f}{{\color{blue} \partial y}}\,\right)
  \]

  \vspace{\baselineskip}
  \[
    f_{{\color{blue} x}{\color{magenta} y}} = \big(\,f_{\color{blue} x}\,\big)_{\color{magenta} y}
  \]
\end{frame}

% 01
%------------------------------------------------------------------------------%
\addtocounter{exemplo}{1}
\begin{frame}{Exemplo \theexemplo}
  Encontre as derivadas de segunda ordem da função
  \[
    f(x, y) = x\cos(y) + ye^x
  \]

  \pause
  \vspace{\baselineskip}
  Note que $f$ é contínua no plano

  \pause
  \vspace{\baselineskip}
  Precisamos calcular \, $f_{xx}$,\, $f_{yy}$,\, $f_{xy}$\, e\, $f_{yx}$
\end{frame}

%------------------------------------------------------------------------------%
\begin{frame}[t]{Exemplo \theexemplo \; -- \; $f_x$}
\begin{align*}
  \onslide<+->{f_x(x, y)}
  \onslide<+->{& = \D{f}{x} \\[2mm]}
  \onslide<+->{& = \D{}{x}\big(x\cos(y) + ye^x\big) \\[2mm]}
  \onslide<+->{& = \D{}{x}\big(x\cos(y)\big) + \D{}{x}\big(ye^x\big) \\[2mm]}
  \onslide<+->{& = \cos(y)\D{x}{x} + y\D{e^x}{x} \\[2mm]}
  \onslide<+->{& = \cos(y) + ye^x }
  \onslide<+->{& \text{(contínua no plano)}}
\end{align*}
\end{frame}

%------------------------------------------------------------------------------%
\begin{frame}[t]{Exemplo \theexemplo \; -- \; $f_y$}
\begin{align*}
  \onslide<+->{f_y(x, y)}
  \onslide<+->{& = \D{f}{y} \\[2mm]}
  \onslide<+->{& = \D{}{y}\big(x\cos(y) + ye^x\big) \\[2mm]}
  \onslide<+->{& = \D{}{y}\big(x\cos(y)\big) + \D{}{y}\big(ye^x\big) \\[2mm]}
  \onslide<+->{& = x\D{}{y}\cos(y) + e^x\D{y}{y} \\[2mm]}
  \onslide<+->{& = -x\sin(y) + e^x \\[2mm]}
  \onslide<+->{& = e^x - x\sin(y)}
  \onslide<+->{& \text{(contínua no plano)}}
\end{align*}
\end{frame}

%------------------------------------------------------------------------------%
\begin{frame}[t]{Exemplo \theexemplo \; -- \; $f_{xx}$}
\begin{align*}
  \onslide<+->{f_{xx}(x, y)}
  \onslide<+->{& = \DS{f}{x} \\[2mm]}
  \onslide<+->{& = \D{}{x}\left(\D{f}{x}\right) \\[2mm]}
  \onslide<+->{& = \D{}{x}\big(\cos(y) + ye^x\big) \\[2mm]}
  \onslide<+->{& = \D{}{x}\big(\cos(y)\big) + \D{}{x}\big(ye^x\big) \\[2mm]}
  \onslide<+->{& = 0 + y\D{e^x}{x} \\[2mm]}
  \onslide<+->{& = ye^x}
  \onslide<+->{& \text{(contínua no plano)}}
\end{align*}
\end{frame}

%------------------------------------------------------------------------------%
\begin{frame}[t]{Exemplo \theexemplo \; -- \; $f_{xy}$}
\begin{align*}
  \onslide<+->{f_{xy}(x, y)}
  \onslide<+->{& = \left(f_x\right)_y(x, y) \\[1mm]}
  \onslide<+->{& = \D{}{y}\left(\D{f}{x}\right) \\[1mm]}
  \onslide<+->{& = \D{}{y}\big(\cos(y) + ye^x\big) \\[1mm]}
  \onslide<+->{& = \D{}{y}\big(\cos(y)\big) + \D{}{y}\big(ye^x\big) \\[1mm]}
  \onslide<+->{& = -\sin(y) + e^x \D{y}{y} \\[1mm]}
  \onslide<+->{& = -\sin(y) + e^x}
  \onslide<+->{  \;=\; e^x - \sin(y)}
  \onslide<+->{& \text{(contínua no plano)}}
\end{align*}
\end{frame}

%------------------------------------------------------------------------------%
\begin{frame}[t]{Exemplo \theexemplo \; -- \; $f_{yy}$}
\begin{align*}
  \onslide<+->{f_{yy}(x, y)}
  \onslide<+->{& = \DS{f}{y} \\[2mm]}
  \onslide<+->{& = \D{}{y}\left(\D{f}{y}\right) \\[1mm]}
  \onslide<+->{& = \D{}{y}\big(e^x - x\sin(y)\big) \\[1mm]}
  \onslide<+->{& = \D{}{y}\big(e^x\big) - \D{}{y}\big(x\sin(y)\big) \\[1mm]}
  \onslide<+->{& = 0 - x\D{}{y}\sin(y) \\[1mm]}
  \onslide<+->{& = -x\cos(x)}
  \onslide<+->{& \text{(contínua no plano)}}
\end{align*}
\end{frame}

%------------------------------------------------------------------------------%
\begin{frame}[t]{Exemplo \theexemplo \; -- \; $f_{yx}$}
\begin{align*}
  \onslide<+->{f_{yx}(x, y) }
  \onslide<+->{& = \left(f_y\right)_x(x, y) \\[1mm]}
  \onslide<+->{& = \D{}{x}\left(\D{f}{y}\right) \\[1mm]}
  \onslide<+->{& = \D{}{x}\big(e^x - x\sin(y)\big) \\[1mm]}
  \onslide<+->{& = \D{}{x}e^x - \D{}{x}\big(x\sin(y)\big) \\[1mm]}
  \onslide<+->{& = e^x - \sin(y)\D{x}{x} \\[1mm]}
  \onslide<+->{& = e^x - \sin(y)}
  \onslide<+->{& \text{(contínua no plano)}}
\end{align*}
\end{frame}

%------------------------------------------------------------------------------%
\begin{frame}{Exemplo \theexemplo \; -- \; Resposta}
\begin{align*}
  f_{xx}(x, y) & = ye^x                           \\[6mm]
  f_{xy}(x, y) & = \color<2>{blue}{e^x - \sin(y)} \\[6mm]
  f_{yy}(x, y) & = -x\cos(x)                      \\[6mm]
  f_{yx}(x, y) & = \color<2>{blue}{e^x - \sin(y)}
\end{align*}
\end{frame}

% 02
%------------------------------------------------------------------------------%
\addtocounter{exemplo}{1}
\begin{frame}{Exemplo \theexemplo}
  Encontre as derivadas de segunda ordem mistas da função
  \(
    f(x, y) = \arctan\left(\frac{y}{x}\right)
  \)

  \pause
  Dica
  \(
    \frac{d}{dt} \arctan(t)
    = \arctan'(t)
    = \frac{1}{1+t^2}
  \)

  \pause
  \vspace{\baselineskip}
  $f$ \emph{não} é contínua no plano

  \pause
  \vspace{\baselineskip}
  Precisamos calcular $f_{xy}$\, e\, $f_{yx}$
\end{frame}

%------------------------------------------------------------------------------%
\begin{frame}{Exemplo \theexemplo \; -- \; $f_x$}
\begin{columns}
\column{0.5\linewidth}
  \begin{align*}
    \onslide<+->{f_x(x, y)}
    \onslide<+->{& = \D{f}{x} \\[2mm]}
    \onslide<+->{& = \D{}{x}\left(\arctan\left(\frac{y}{x}\right)\right) \\[2mm]}
    \onslide<+->{& = \arctan'\left(\frac{y}{x}\right)\D{}{x}\left(\frac{y}{x}\right) \\[2mm]}
    \onslide<+->{& = \arctan'\left(\frac{y}{x}\right)y\D{x^{-1}}{x} \\[2mm]}
    \onslide<+->{& = \arctan'\left(\frac{y}{x}\right)y\left(-x^{-2}\right) \\[2mm]}
    \onslide<+->{& = \arctan'\left(\frac{y}{x}\right)\frac{-y}{x^2}}
  \end{align*}
\column{0.5\linewidth}
  \onslide<+->{Lembrando que \( \arctan'(t) = \frac{1}{1+t^2} \)}
  \begin{align*}
    \onslide<+->{f_x(x, y)
    & = \frac{1}{1+\left(\dfrac{y}{x}\right)^2}     \;\frac{-y}{x^2} \\[2mm]}
    \onslide<+->{& = \frac{1}{\dfrac{x^2+y^2}{x^2}} \;\frac{-y}{x^2} \\[2mm]}
    \onslide<+->{& = \frac{x^2}{x^2+y^2} \;\frac{-y}{x^2} \\[2mm]}
    \onslide<+->{& = \frac{-y}{x^2+y^2}}
  \end{align*}
\end{columns}
\end{frame}

%------------------------------------------------------------------------------%
\begin{frame}{Exemplo \theexemplo \; -- \; $f_y$}
\begin{columns}
  \column{0.5\linewidth}
  \begin{align*}
    \onslide<+->{f_y(x, y)}
    \onslide<+->{& = \D{f}{y} \\[2mm]}
    \onslide<+->{& = \D{}{y}\left(\arctan\left(\frac{y}{x}\right)\right)             \\[2mm]}
    \onslide<+->{& = \arctan'\left(\frac{y}{x}\right)\D{}{y}\left(\frac{y}{x}\right) \\[2mm]}
    \onslide<+->{& = \arctan'\left(\frac{y}{x}\right)\frac{1}{x}\D{y}{y}             \\[2mm]}
    \onslide<+->{& = \arctan'\left(\frac{y}{x}\right)\frac{1}{x}}
  \end{align*}
  \column{0.5\linewidth}
  \onslide<+->{Lembrando que \( \arctan'(t) = \frac{1}{1+t^2} \)}
  \begin{align*}
    \onslide<+->{f_y(x, y)
      & = \frac{1}{1+\left(\dfrac{y}{x}\right)^2}   \;\frac{1}{x} \\[1mm]}
    \onslide<+->{& = \frac{1}{\dfrac{x^2+y^2}{x^2}} \;\frac{1}{x} \\[1mm]}
    \onslide<+->{& = \frac{x^2}{x^2+y^2}\;\frac{1}{x}             \\[1mm]}
    \onslide<+->{& = \frac{x}{x^2+y^2}}
  \end{align*}
\end{columns}
\end{frame}

%------------------------------------------------------------------------------%
\begin{frame}[t]{Exemplo \theexemplo \; -- \; $f_{xy}$}
\begin{columns}[t]
\column{0.4\linewidth}
    \begin{align*}
    \onslide<+->{f_{xy}(x, y)}
    \onslide<+->{& = \left(f_x\right)_y(x, y) \\[1mm]}
    \onslide<+->{& = \D{}{y}\left(\D{f}{x}\right) \\[1mm]}
    \onslide<+->{& = \D{}{y}\left(\frac{-y}{x^2+y^2}\right) \\[1mm]}
    \onslide<+->{& = -\frac{\displaystyle (x^2+y^2)\D{y}{y}-y\D{}{y}(x^2+y^2)}{(x^2+y^2)^2}}
  \end{align*}
\column{0.6\linewidth}
  \begin{align*}
    \onslide<+->{f_{xy}(x, y)
    & = -\frac{\displaystyle x^2+y^2-y\left(\D{x^2}{y}+\D{y^2}{y}\right)}{(x^2+y^2)^2} \\[1mm]}
    \onslide<+->{& = -\frac{x^2+y^2-y(2y)}{(x^2+y^2)^2} \\[2mm]}
    \onslide<+->{& = -\frac{x^2-y^2}{(x^2+y^2)^2} \\[2mm]}
    \onslide<+->{& = \frac{y^2-x^2}{(x^2+y^2)^2}}
  \end{align*}
\end{columns}
\end{frame}

%------------------------------------------------------------------------------%
\begin{frame}[t]{Exemplo \theexemplo \; -- \; $f_{yx}$}
\begin{columns}
\column{0.4\linewidth}
  \begin{align*}
    \onslide<+->{f_{yx}(x, y)}
    \onslide<+->{& = \left(f_y\right)_x(x, y) \\[1mm]}
    \onslide<+->{& = \D{}{x}\left(\D{f}{y}\right) \\[1mm]}
    \onslide<+->{& = \D{}{x}\left(\frac{x}{x^2+y^2}\right) \\[1mm]}
    \onslide<+->{& = \frac{\displaystyle \left(x^2+y^2\right)\D{x}{x} - x \D{}{x}\left(x^2+y^2\right)}
             {\left(x^2+y^2\right)^2}}
  \end{align*}
\column{0.6\linewidth}
  \begin{align*}
    \onslide<+->{f_{yx}(x, y)
    & = \frac{\displaystyle x^2+y^2 - x \left(\D{x^2}{x}+\D{y^2}{x}\right)}
             {\left(x^2+y^2\right)^2} \\[2mm]}
    \onslide<+->{& = \frac{x^2+y^2 - x \left(2x\right)}{\left(x^2+y^2\right)^2} \\[2mm]}
    \onslide<+->{& = \frac{y^2 - x^2}{\left(x^2+y^2\right)^2}}
  \end{align*}
\end{columns}
\end{frame}

%------------------------------------------------------------------------------%
\begin{frame}{Exemplo \theexemplo \; -- \; Resposta}
  \begin{align*}
    f_{xy}(x, y) & = \frac{y^2-x^2}{\left(x^2+y^2\right)^2} \\[6mm]
    f_{yx}(x, y) & = \frac{y^2-x^2}{\left(x^2+y^2\right)^2}
  \end{align*}
\end{frame}

%------------------------------------------------------------------------------%
\section{Teorema das Derivadas Mistas}

%------------------------------------------------------------------------------%
\begin{frame}{Teorema das Derivadas Mistas}
  Se a função \(f(x, y)\) e suas derivadas parciais
  \, $f_\emph{x}$,\, $f_\emph{y}$,\, $f_{\emph{xy}}$\, e \,$f_{\emph{yx}}$

  \pause
  forem definidas em uma vizinhança do ponto \((a, b)\)

  \pause
  e \emph{todas forem contínuas}, então

  \pause
  \[
    \DM{f}{x}{y} = \DM{f}{y}{x}
  \]
\end{frame}

%------------------------------------------------------------------------------%
\addtocounter{exemplo}{1}
\begin{frame}{Exemplo \theexemplo}
  Encontre
  \(\;
    \DM{w}{x}{y}
  \;\)
  se
  \(\;
    w = xy + \frac{e^y}{y^2+1}
  \;\)
\end{frame}

%------------------------------------------------------------------------------%
\begin{frame}{Exemplo \theexemplo}
  \onslide<+->{Queremos calcular
  \[
    \DM{w}{x}{y} = \D{}{x}\left(\D{w}{y}\right)
  \]}
  \onslide<+->{Calculando a primeira derivada}
  \begin{align*}
    \onslide<+->{\D{w}{y}
    & = \D{}{y}\left(xy + \frac{e^y}{y^2+1}\right) \\[2mm]}
    \onslide<+->{& = \D{}{y}\left(xy\right) + \D{}{y}\left(\frac{e^y}{y^2+1}\right)}
  \end{align*}

  \onslide<+->{Precisamos calcular a derivada de uma divisão}
\end{frame}

%------------------------------------------------------------------------------%
\begin{frame}{Exemplo \theexemplo}
  Observando que $w$ é contínua no plano

  \pause
  e \emph{assumindo} que suas derivadas parciais também serão

  \pause
  podemos usar o teorema
  \[
    \DM{w}{x}{y} = \DM{w}{y}{x}
  \]

  \pause
  Temos então que calcular
  \[
    \DM{w}{x}{y} = \D{}{y}\left(\D{w}{x}\right)
  \]
\end{frame}

%------------------------------------------------------------------------------%
\begin{frame}{Exemplo \theexemplo}
  \onslide<+->{Calculando a derivada primeira}
  \begin{align*}
    \onslide<+->{\D{w}{x}
    & = \D{}{x}\left(xy + \frac{e^y}{y^2+1}\right) \\[2mm]}
    \onslide<+->{& = \D{}{x}\left(xy\right) + \D{}{x}\left(\frac{e^y}{y^2+1}\right) \\[2mm]}
    \onslide<+->{& = y \D{x}{x} + 0 \\[2mm]}
    \onslide<+->{& = y}
    \onslide<+->{& \text{contínua no plano}}
  \end{align*}
\end{frame}

%------------------------------------------------------------------------------%
\begin{frame}{Exemplo \theexemplo}
  \onslide<+->{Calculando a derivada segunda}
  \begin{align*}
    \onslide<+->{\DM{w}{x}{y}
    & = \D{}{y}\left(\D{w}{x}\right) \\[2mm]}
    \onslide<+->{& = \D{}{y}\left(y\right) \\[2mm]}
    \onslide<+->{& = 1}
    \onslide<+->{& \text{contínua no plano}}
  \end{align*}
\end{frame}

%------------------------------------------------------------------------------%
\section{Derivadas Parciais de Ordem Superior}

%------------------------------------------------------------------------------%
\begin{frame}{Notações}
  \begin{columns}
    \column{0.5\linewidth}
    Derivadas terceiras
    \[
      \frac{\partial^3 f}{\partial x^3} \qquad\qquad\qquad
      f_{xxx}
    \]

    \vspace{\baselineskip}
    \[
      \frac{\partial^3 f}{\partial y\partial x^2} \qquad\qquad\qquad
      f_{xxy}
    \]
    \column{0.5\linewidth}
    Derivadas quartas
    \[
      \frac{\partial^4 f}{\partial x^4} \qquad\qquad\qquad
      f_{xxx}
    \]

    \vspace{\baselineskip}
    \[
      \frac{\partial^4 f}{\partial y^2 \partial x^2} \qquad\qquad\qquad
      f_{xxyy}
    \]
  \end{columns}
\end{frame}

%------------------------------------------------------------------------------%
\begin{frame}{Ordem de Derivação}
  Desde que todas as derivadas envolvidas sejam contínuas

  \pause
  a ordem de derivação é irrelevante
\end{frame}

%------------------------------------------------------------------------------%
\addtocounter{exemplo}{1}
\begin{frame}{Exemplo \theexemplo}
  Encontre
  \(
    f_{yxyz}
  \)
  da função
  \(
    f(x, y, z) = 1 - 2xy^2z + x^2y
  \)

  \pause
  \vspace{\baselineskip}
  Queremos calcular
  \[
    f_{yxyz} = \big(\,\big(\,\big(\,f_y\,\big)\!\,_x\,\big)\!\,_y\,\big)\!\,_z
  \]
\end{frame}

%------------------------------------------------------------------------------%
\begin{frame}{Exemplo \theexemplo}
\begin{align*}
  \onslide<+->{f        & = 1 - 2xy^2z + x^2y} \\
  \onslide<+->{f_y      & = \D{}{y}\left(1 - 2xy^2z + x^2y\right)}
  \onslide<+->{           = 0 -2x2yz + x^2}
  \onslide<+->{           = -4xyz + x^2 \\}
  \onslide<+->{f_{yx}   & = \D{}{x}\left(-4xyz + x^2\right)}
  \onslide<+->{           = -4yz + 2x \\}
  \onslide<+->{f_{yxy}  & = \D{}{y}\left(-4yz + 2x\right)}
  \onslide<+->{           = -4z + 0}
  \onslide<+->{           = -4z \\}
  \onslide<+->{f_{yxyz} & = \D{}{z}\left(-4z\right)}
  \onslide<+->{           = -4}
\end{align*}
\end{frame}

%------------------------------------------------------------------------------%
\section{Lista Mínima}

%------------------------------------------------------------------------------%
\begin{frame}{Lista Mínima}
  Cálculo Vol. 2 do Thomas 12\textsuperscript{a} ed. -- Seção 14.3
  \begin{enumerate}
    \item Estudar o texto da seção
    \item Resolver os exercícios: 43-45, 51-54, 57, 61
  \end{enumerate}

  \vfill
  Atenção: A prova é baseada no livro, não nas apresentações
\end{frame}

%------------------------------------------------------------------------------%
\end{document}
%------------------------------------------------------------------------------%
